\documentclass[a4paper]{article}
\usepackage{amsfonts}
\usepackage{amsmath}
\usepackage{amssymb}

\begin{document}

\noindent \large Auteur: Dina Haajou \\
\noindent \large Studierichting: Informatica \\
\noindent \large Oefeningenreeks 1E nr. 3 \\

\normalsize

We veronderstellen dat $f(z)$ gedeeld wordt door $(z + 1)(z - 1) = z^2 - 1$. \\

De rest is dan van de vorm $az + b$. \\

$f(1) = 5$ \\

$f(-1) = 1$ \\

$A(z) = (z - 1)(z + 1) + az + b$ \\

$A(1) = a \cdot 1 + b = 5$ \\

$A(-1) = a \cdot (-1) + b = 1$ \\

$\Rightarrow \left\{
   \begin{array}{l}
     a + b = 5 \\
     -a + b = 1 \\
   \end{array}
\right.$ \\

$\Rightarrow \left\{
   \begin{array}{l}
     a = 2 \\
     b = 3 \\
   \end{array}
\right.$ \\

$(a, b) = (2, 3)$ \\

R = $2z + 3$

\begin{thebibliography}{999}
%\bibitem{a} Referentie
\end{thebibliography}

\end{document}
